% GENERATED FILE. Edit canonical lessons, then run npm run generate:books.
% canonical-source-hash: fnv1a64:fc6e75330c4917b4
% canonical-lessons: TE-C425-nirnayam, TE-C425-tirmanam, TE-C425-telcu, TE-C425-khararu-ceyu, TE-C425-punaralocincu

\chapter{Making a Decision}
\label{ch:te-making-a-decision}
\hlchaptermodality{425}

\begin{chapteropening}
\textbf{By the end of this chapter:} \emph{I can talk about decisions, resolutions and rethinking.}

Complete the last lesson of chapter 425: I can talk about decisions, resolutions and rethinking.
\end{chapteropening}

\section[\texorpdfstring{\te{నిర్ణయం} (nirṇayaṁ)}{నిర్ణయం (nirṇayaṁ)}]{\te{నిర్ణయం} (nirṇayaṁ) — a decision}

\label{lesson:TE-C425-nirnayam}



\begin{quote}
Before the new one: say the Telugu for to stand up to, then the Telugu for to withdraw.
\end{quote}

\subsection*{You'll want to know: \te{నిర్ణయం}}
\textbf{\te{నిర్ణయం}} — \emph{nirṇayaṁ} — \textquotedblleft{}a decision\textquotedblright{}.

\begin{quote}
\te{ఆమె} \te{మంచి} \te{నిర్ణయం} \te{తీసుకుంది}.
\end{quote}

\begin{quote}
\emph{āme mañci nirṇayaṁ tīsukundi.}
\end{quote}

\begin{quote}
She made a good decision.
\end{quote}

Plural: \te{నిర్ణయాలు}, \emph{nirṇayālu}.

\te{నిర్ణయం} \te{తీసుకో}, to take a decision; \te{నిర్ణయించు}, to decide, which you know.

\begin{cousinweb}
\textbf{\te{నిర్ణయం}} (\emph{nirṇayaṁ}) is borrowed from Sanskrit \emph{nirṇaya}, \textquotedblleft{}complete ascertainment, decision, settlement\textquotedblright{} (Monier-Williams); \te{నిర్ణయించు} is its verb.
\end{cousinweb}

\begin{grammarlens}[title={What you've built}]
One more word for saying what happens if.
\end{grammarlens}

\subsection*{Your turn}
\begin{itemize}
\raggedright
  \item \emph{Say it:} \emph{nirṇayaṁ}
  \item \emph{Say it:} \emph{nirṇayaṁ}, once more
  \item \emph{Say it:} \emph{āme mañci nirṇayaṁ tīsukundi.} — She made a good decision.
  \item \emph{Recall:} say the Telugu for to stand up to, then the Telugu for to withdraw, then say \emph{nirṇayaṁ} again
\end{itemize}

\begin{tcolorbox}[breakable,colback=teal!4,colframe=teal!35!black,title={Before you move on}]
What does \te{నిర్ణయం} mean? (\textquotedblleft{}A decision\textquotedblright{}.) Say it once more.
\end{tcolorbox}
\section[\texorpdfstring{\te{తీర్మానం} (tīrmānaṁ)}{తీర్మానం (tīrmānaṁ)}]{\te{తీర్మానం} (tīrmānaṁ) — a resolution (of a meeting)}

\label{lesson:TE-C425-tirmanam}



\begin{quote}
Before the new one: say the Telugu for to withdraw, then the Telugu for a decision.
\end{quote}

\subsection*{You'll want to know: \te{తీర్మానం}}
\textbf{\te{తీర్మానం}} — \emph{tīrmānaṁ} — \textquotedblleft{}a resolution (of a meeting)\textquotedblright{}.

\begin{quote}
\te{సమావేశంలో} \te{ఒక} \te{తీర్మానం} \te{ఆమోదించారు}.
\end{quote}

\begin{quote}
\emph{samāvēśaṁlō oka tīrmānaṁ āmōdiñcāru.}
\end{quote}

\begin{quote}
The meeting passed a resolution.
\end{quote}

Plural: \te{తీర్మానాలు}, \emph{tīrmānālu}.

\te{తీర్మానం} \te{చేయు}, to resolve formally.

\begin{cousinweb}
\textbf{\te{తీర్మానం}} (\emph{tīrmānaṁ}) is native Dravidian: DEDR 3278 lists Kannada \emph{tīrmāna}, \textquotedblleft{}conclusion, settlement, decision\textquotedblright{}, beside Telugu \emph{tīru}, \textquotedblleft{}to be finished, to be settled\textquotedblright{}.
\end{cousinweb}

\begin{grammarlens}[title={What you've built}]
One more word for saying what happens if.
\end{grammarlens}

\subsection*{Your turn}
\begin{itemize}
\raggedright
  \item \emph{Say it:} \emph{tīrmānaṁ}
  \item \emph{Say it:} \emph{tīrmānaṁ}, once more
  \item \emph{Say it:} \emph{samāvēśaṁlō oka tīrmānaṁ āmōdiñcāru.} — The meeting passed a resolution.
  \item \emph{Recall:} say the Telugu for to withdraw, then the Telugu for a decision, then say \emph{tīrmānaṁ} again
\end{itemize}

\begin{tcolorbox}[breakable,colback=teal!4,colframe=teal!35!black,title={Before you move on}]
What does \te{తీర్మానం} mean? (\textquotedblleft{}A resolution (of a meeting)\textquotedblright{}.) Say it once more.
\end{tcolorbox}
\section[\texorpdfstring{\te{తేల్చు} (tēlcu)}{తేల్చు (tēlcu)}]{\te{తేల్చు} (tēlcu) — to sort out; to settle}

\label{lesson:TE-C425-telcu}



\begin{quote}
Before the new one: say the Telugu for a decision, then the Telugu for a resolution.
\end{quote}

\subsection*{You'll want to know: \te{తేల్చు}}
\textbf{\te{తేల్చు}} — \emph{tēlcu} — \textquotedblleft{}to sort out; to settle\textquotedblright{}.

\begin{quote}
\te{ఆ} \te{గొడవను} \te{పెద్దలు} \te{తేల్చారు}.
\end{quote}

\begin{quote}
\emph{ā goḍavanu peddalu tēlcāru.}
\end{quote}

\begin{quote}
The elders settled that quarrel.
\end{quote}

Past (I …): \te{తేల్చాను}, \emph{tēlcānu}.

\te{తేల్చి} \te{చెప్పు}, to say plainly, once and for all.

\begin{cousinweb}
\textbf{\te{తేల్చు}} (\emph{tēlcu}) is native Dravidian: DEDR 3464 lists Telugu \emph{tēlucu}, \textquotedblleft{}to cause to float\textquotedblright{}, beside \emph{tēlu}, \textquotedblleft{}to float\textquotedblright{}; in everyday Telugu it means to settle a matter.
\end{cousinweb}

\begin{grammarlens}[title={What you've built}]
One more word for saying what happens if.
\end{grammarlens}

\subsection*{Your turn}
\begin{itemize}
\raggedright
  \item \emph{Say it:} \emph{tēlcu}
  \item \emph{Say it:} \emph{tēlcu}, once more
  \item \emph{Say it:} \emph{ā goḍavanu peddalu tēlcāru.} — The elders settled that quarrel.
  \item \emph{Recall:} say the Telugu for a decision, then the Telugu for a resolution, then say \emph{tēlcu} again
\end{itemize}

\begin{tcolorbox}[breakable,colback=teal!4,colframe=teal!35!black,title={Before you move on}]
What does \te{తేల్చు} mean? (\textquotedblleft{}To sort out; to settle\textquotedblright{}.) Say it once more.
\end{tcolorbox}
\section[\texorpdfstring{\te{ఖరారు} \te{చేయు} (kharāru cēyu)}{ఖరారు చేయు (kharāru cēyu)}]{\te{ఖరారు} \te{చేయు} (kharāru cēyu) — to finalise; to fix (a date)}

\label{lesson:TE-C425-khararu-ceyu}



\begin{quote}
Before the new one: say the Telugu for a resolution, then the Telugu for to sort out.
\end{quote}

\subsection*{You'll want to know: \te{ఖరారు} \te{చేయు}}
\textbf{\te{ఖరారు} \te{చేయు}} — \emph{kharāru cēyu} — \textquotedblleft{}to finalise; to fix (a date)\textquotedblright{}.

\begin{quote}
\te{పెళ్ళి} \te{తేదీని} \te{ఖరారు} \te{చేశారు}.
\end{quote}

\begin{quote}
\emph{peḷḷi tēdīni kharāru cēśāru.}
\end{quote}

\begin{quote}
They fixed the date of the wedding.
\end{quote}

Past (I …): \te{ఖరారు} \te{చేశాను}, \emph{kharāru cēśānu}.

\te{ఖరారు} alone: confirmation, a settled fact.

\begin{cousinweb}
\textbf{\te{ఖరారు} \te{చేయు}} (\emph{kharāru cēyu}) is a loan from Arabic \emph{qarār}, \textquotedblleft{}settlement, a fixed state\textquotedblright{}, with the native verb \te{చేయు}, \textquotedblleft{}to do\textquotedblright{} (DEDR 1957).
\end{cousinweb}

\begin{grammarlens}[title={What you've built}]
One more word for saying what happens if.
\end{grammarlens}

\subsection*{Your turn}
\begin{itemize}
\raggedright
  \item \emph{Say it:} \emph{kharāru cēyu}
  \item \emph{Say it:} \emph{kharāru cēyu}, once more
  \item \emph{Say it:} \emph{peḷḷi tēdīni kharāru cēśāru.} — They fixed the date of the wedding.
  \item \emph{Recall:} say the Telugu for a resolution, then the Telugu for to sort out, then say \emph{kharāru cēyu} again
\end{itemize}

\begin{tcolorbox}[breakable,colback=teal!4,colframe=teal!35!black,title={Before you move on}]
What does \te{ఖరారు} \te{చేయు} mean? (\textquotedblleft{}To finalise; to fix (a date)\textquotedblright{}.) Say it once more.
\end{tcolorbox}
\section[\texorpdfstring{\te{పునరాలోచించు} (punarālōciñcu)}{పునరాలోచించు (punarālōciñcu)}]{\te{పునరాలోచించు} (punarālōciñcu) — to reconsider; to rethink}

\label{lesson:TE-C425-punaralocincu}



\begin{quote}
Before the new one: say the Telugu for to sort out, then the Telugu for to finalise.
\end{quote}

\subsection*{You'll want to know: \te{పునరాలోచించు}}
\textbf{\te{పునరాలోచించు}} — \emph{punarālōciñcu} — \textquotedblleft{}to reconsider; to rethink\textquotedblright{}.

\begin{quote}
\te{అతను} \te{తన} \te{నిర్ణయాన్ని} \te{పునరాలోచించాడు}.
\end{quote}

\begin{quote}
\emph{atanu tana nirṇayānni punarālōciñcāḍu.}
\end{quote}

\begin{quote}
He reconsidered his decision.
\end{quote}

Past (I …): \te{పునరాలోచించాను}, \emph{punarālōciñcānu}.

\te{పునర్}-, again, + \te{ఆలోచించు}, to think over.

\begin{cousinweb}
\textbf{\te{పునరాలోచించు}} (\emph{punarālōciñcu}) is built from Sanskrit \emph{punar}, \textquotedblleft{}back, again\textquotedblright{} (CDIAL 8273), and \emph{ālocana}, \textquotedblleft{}considering, reflection\textquotedblright{} (Monier-Williams), with the Telugu verb ending -\te{ఇంచు}.
\end{cousinweb}

\begin{grammarlens}[title={What you've built}]
One more word for saying what happens if.
\end{grammarlens}

\subsection*{Your turn}
\begin{itemize}
\raggedright
  \item \emph{Say it:} \emph{punarālōciñcu}
  \item \emph{Say it:} \emph{punarālōciñcu}, once more
  \item \emph{Say it:} \emph{atanu tana nirṇayānni punarālōciñcāḍu.} — He reconsidered his decision.
  \item \emph{Recall:} say the Telugu for to sort out, then the Telugu for to finalise, then say \emph{punarālōciñcu} again
\end{itemize}

\begin{tcolorbox}[breakable,colback=teal!4,colframe=teal!35!black,title={Before you move on}]
What does \te{పునరాలోచించు} mean? (\textquotedblleft{}To reconsider; to rethink\textquotedblright{}.) Say it once more.
\end{tcolorbox}
